\documentclass[../main.tex]{subfiles}
\begin{document}
\section{Pentesting}

El \textbf{pentesting} (penetration testing o prueba de penetración) es una de las
formas de corroborar el funcionamiento correcto de un sistema en cuanto a seguridad.
Parte de la premisa de que \emph{ya existen vulnerabilidades} y trata de probar que
esas vulnerabilidades existen.

\begin{destacado}
Esta unidad corresponde al cierre del temario y \textbf{entra completa en el parcial}.
El profe lo remarcó: \prof{Esto no es así un corolario, sino que va a entrar en el parcial.}
Lo más importante a saber de memoria es la \textbf{metodología de hipótesis de falla}
(los pasos en orden) y la idea de que el pentesting \textbf{no garantiza seguridad}.
\end{destacado}

\subsection{Formas de verificar un sistema: verificación formal vs.\ pentesting}

Hay varias formas de corroborar el funcionamiento correcto de un sistema. Dos de ellas:

\begin{definicion}{Verificación formal}
Manera de comprobar \textbf{matemáticamente} que un sistema cumple con las restricciones
que le imponemos. Esquema:
\begin{itemize}
    \item \textbf{Precondiciones} $\rightarrow$ hipótesis sobre el estado del sistema
    (por ejemplo: el sistema está en un estado seguro).
    \item Se ejecutan las operaciones del sistema ($Op_1 \dots Op_n$) sobre un input dado.
    \item \textbf{Poscondiciones} $\rightarrow$ resultado de aplicar esas operaciones al input.
    \item \textbf{Requerimiento}: las poscondiciones cumplen las restricciones pedidas.
\end{itemize}
Sirve para probar matemáticamente que un sistema hace lo que dice hacer.
\end{definicion}

Características de la verificación formal según el profe:
\begin{itemize}
    \item Es \textbf{súper compleja}: con un par de variables y operaciones ya es difícil
    de seguir; hay que mapear todas las posibilidades, precondiciones e interacciones
    entre componentes.
    \item Es \textbf{muy costosa y extensa}; a veces termina siendo \textbf{impráctica}:
    puede llevar más tiempo hacer la verificación formal que hacer el propio sistema.
    \item Peor aún cuando el sistema \textbf{evoluciona rápido}: antes de terminar la
    verificación ya quedó obsoleta y hay que hacer otra.
    \item Es realizable a nivel de funciones (como las viejas pruebas de escritorio:
    pruebo un input, otro input, qué pasa si le paso \texttt{null}), pero a nivel de un
    sistema completo se vuelve poco práctica.
\end{itemize}

\begin{definicion}{Prueba de penetración (pentesting)}
\emph{Prueba} para verificar que un sistema cumple con ciertas restricciones.
\begin{itemize}
    \item \textbf{Precondiciones} $\rightarrow$ hipótesis sobre el estado del sistema
    \textbf{y la existencia de una vulnerabilidad}.
    \item \textbf{Poscondiciones} $\rightarrow$ sistema comprometido.
    \item \textbf{Ejecución}: aplicar pruebas para intentar mover al sistema del estado
    inicial al estado comprometido.
\end{itemize}
Es mucho más \textbf{económica, simple y eficiente} que la verificación formal, y es lo
que hacen las empresas. Permite probar primero las áreas/componentes más críticos y
después los menos críticos.
\end{definicion}

\subsubsection*{Similitudes y diferencias}

\begin{itemize}
    \item \textbf{Ambas} tratan de probar la \textbf{existencia} de vulnerabilidades.
    \item La \textbf{prueba de penetración NO prueba la ausencia} de vulnerabilidades.
    \item La \textbf{verificación formal SÍ prueba la ausencia}, pero para ello debe
    incluir \textbf{TODOS} los factores posibles (factores externos).
    \item En la práctica, la verificación formal prueba ausencia de vulnerabilidades en
    determinado algoritmo, programa o ambiente, \textbf{no en un sistema completo}. Se
    ignoran cosas como la \textbf{instalación} y el \textbf{uso}, el sistema operativo
    donde se instala, etc.
\end{itemize}

\begin{destacado}
El profe lo conecta con Dijkstra: el testing \textbf{puede mostrar la presencia} de bugs
o vulnerabilidades, \textbf{pero no prueba su ausencia}. Esto va en total concordancia
con el pentesting.
\end{destacado}

\begin{ojo}
Que la poscondición sea ``no encontré nada'' / ``no se dio la hipótesis'' \textbf{no
significa} que no haya alguna falla. Solo significa que el sistema \textbf{no está
comprometido bajo las condiciones de esa prueba}. Puede haber otros vectores de ataque
que no se probaron.
\end{ojo}

\subsection{Objetivos y marco ético (ethical hacking)}

El objetivo es \textbf{probar la eficacia de los controles de seguridad de un sistema},
intentando \textbf{violar la política de seguridad} existente. Requiere ejecutar técnicas
similares a las de un atacante; por eso se lo llama \textbf{\emph{ethical hacking}}.

\begin{itemize}
    \item Generalmente hay un \textbf{equipo designado} que sabe de seguridad y conoce las
    implicancias. Hay que \textbf{simular ser un atacante} pero adoptando un comportamiento
    ético: no robar datos de clientes, no hacer compras a nombre de otra persona, etc.
    \item Es deseable que lo haga \textbf{alguien que no haya intervenido en la construcción}
    del sistema (análogo a QA: el desarrollador hace pruebas unitarias / de iteración /
    regresiones, y después un QA trata de romperlo con inputs inválidos, clics inesperados,
    etc.).
    \item Es \textbf{análogo a pruebas manuales} del sistema y \textbf{prueba al sistema como
    un todo}, no solo aspectos técnicos.
    \item \textbf{NO reemplaza un buen diseño e implementación}: estos tienen que estar de
    base y ser un requerimiento del sistema.
\end{itemize}

\subsection{Metodología informal (preparación)}

Vista como la etapa de preparación previa:
\begin{itemize}
    \item \textbf{Determinar y cuantificar objetivos} (ej.: obtener información de clientes).
    \item Encontrar cierta cantidad de vulnerabilidades, o buscarlas durante un período de
    tiempo. El estudio se conduce \textbf{desde el punto de vista de un atacante}.
    \item Definir alcance, tiempo, propósito; armar el \textbf{plan de ataque}; elegir
    herramientas; seleccionar el equipo e identificar a quiénes van a participar y a estar
    al tanto (a veces \textbf{no se avisa}; lo del \emph{lower line} / aviso a desarrolladores
    y equipos operativos se decide acá).
    \item Definir una condición de \textbf{objetivo} (éxito/no éxito). El ataque termina al
    lograr el objetivo o al vencerse el tiempo definido.
    \item \textbf{Estudiar y categorizar los hallazgos}: el éxito de una prueba de penetración
    proviene del \textbf{análisis de los hallazgos}. Es un proceso \textbf{cíclico}: los
    hallazgos se reincorporan al ciclo, permitiendo detectar problemas recurrentes y corregir
    sistemáticamente familias de problemas (una misma vulnerabilidad puede repetirse en
    varios componentes por usar la misma dependencia o forma de desarrollo).
\end{itemize}

\subsection{Metodología de Hipótesis de Falla (LOS PASOS)}

\begin{destacado}
\prof{Esta metodología es importante, súper importante que la entiendan y la sepan; es la
base de lo que hoy se usa para testear.} Es la base para equipos de Red Team, AppSec e
InfraSec. Para sistemas medianos o grandes hoy se exige tener algún tipo de pentesting,
sobre todo en los componentes \emph{core}, donde se pide pasar el pentesting sin
vulnerabilidades críticas o altas (definidas por los CVE; importancia de la organización
\textbf{MITRE}).
\end{destacado}

Los \textbf{pasos en orden} (tal como aparecen en el PDF):

\begin{enumerate}
    \item \textbf{Recolección de información} --- para conocer el sistema y su funcionamiento.
    \item \textbf{Hipótesis} (\textbf{Planteo de hipótesis de falla}) --- asumir la existencia
    de vulnerabilidades concretas.
    \item \textbf{Prueba} --- probar las vulnerabilidades.
    \item \textbf{Generalización} --- generalizar patrones y hallar otras vulnerabilidades del
    mismo tipo.
    \item \textbf{Eliminación} (\textbf{opcional}) --- determinar los pasos necesarios para
    eliminarlas.
\end{enumerate}

\begin{examen}
\faExclam \textbf{Foco de parcial (2025-2Q): los pasos del pentesting.} El orden correcto
incluye explícitamente el \textbf{``Planteo de hipótesis de falla''}. Secuencia tipo:
\[
\text{Recolección de info} \rightarrow \text{Planteo de hipótesis de falla} \rightarrow
\text{Prueba} \rightarrow \text{Generalización} \rightarrow \text{Eliminación}
\]
El profe lo llama conocimiento ``enciclopédico'': hay que saberlo de memoria.
\textbf{Equivocar el orden o los pasos PENALIZA.} Notar que la \textbf{Eliminación es
opcional} (solo se incluyen recomendaciones).
\end{examen}

\subsubsection{Paso 1: Recolección de información}

\begin{itemize}
    \item \textbf{Componer un modelo del sistema y sus partes}. Hay que conocer bien el
    sistema para saber ``dónde poner el ojo''.
    \item \textbf{Buscar discrepancias}: cuando el sistema evoluciona y la documentación no
    refleja esa evolución (la doc dice que se comporta de una forma y ya no es así)
    $\rightarrow$ posible punto de entrada.
    \item \textbf{Revisar interfaces}: conexiones con sistemas externos, proveedores, otros
    sistemas, bases de datos.
    \item Buscar \textbf{vulnerabilidades conocidas genéricas} compatibles con la tecnología
    que usa el sistema.
    \item Apoyarse en la \textbf{documentación} (cuanta más, mejor); prestar especial atención
    a \textbf{secciones poco especificadas o ambiguas} (donde pudo tomarse una decisión
    cuestionable $\rightarrow$ otra puerta de entrada).
    \item Ver \textbf{manejo de privilegios y tipos de cuentas}: enumerar usuarios, buscar
    cuentas \emph{superadmin} (las que más posibilidades dan, no están limitadas), cuentas
    \textbf{huérfanas} o con \textbf{exceso de permisos}.
    \item Conocer el \textbf{ambiente}: nombres de usuarios/servidores, en qué servidores
    corre, configuración y \textbf{estructura de red}.
\end{itemize}

\paragraph{Punto de partida.} Identifica la información inicial con la que cuenta el equipo
(supuesto atacante). \textbf{Niveles incrementales}:
\begin{enumerate}
    \item Atacante externo \textbf{sin} conocimiento del sistema.
    \item Atacante externo \textbf{con} conocimiento del sistema.
    \item Atacante \textbf{con acceso} al sistema (ya tiene usuario y contraseña).
\end{enumerate}
Dependiendo del tipo de prueba, algunos niveles son irrelevantes: durante el diseño se
ignora el primer nivel; en sistemas con registro abierto se ignora el segundo nivel.

\subsubsection{Paso 2: Hipótesis}

Buscar posibles vulnerabilidades / combinaciones que expongan una vulnerabilidad:
\begin{itemize}
    \item \textbf{Examinando políticas y procedimientos}: buscar inconsistencias, e
    inconsistencias entre políticas y mecanismos. Los procedimientos pueden \textbf{no
    cumplirse}.
    \item \textbf{Examinando implementaciones}: usar modelos de vulnerabilidades; revisar
    vulnerabilidades comunes/conocidas (\emph{vulnerability scanning}, ej.\ \texttt{portmapper});
    usar manuales (exceder límites, omitir pasos de las secuencias, etc.).
    \item \textbf{Examinando mecanismos}: pueden estar mal implementados, el ambiente donde
    corren puede introducir errores, o pueden no ser seguros.
    \item \textbf{Comparando con otros sistemas}: sistemas parecidos tienen problemas parecidos.
\end{itemize}

Áreas concretas que el profe menciona como candidatas de hipótesis:
\begin{itemize}
    \item \textbf{Seguridad en APIs}: política de desarrollo de \textbf{microservicios}
    (más servicios = más ``cajitas'' atacables; si cada equipo no baja la misma política de
    seguridad, hay diseño menos robusto). \textbf{Software Supply Chain (API) Management}:
    hoy es \textbf{top 3} del OWASP Top Ten para apps web; abarca dependencias vulnerables,
    infiltración en flujos CI/CD con plugins/dependencias vulnerables, y ataques de
    confusión de paquetes (ej.\ casos de \texttt{npm}, donde se baja todo el árbol de
    dependencias transitivas).
    \item \textbf{RASP (Runtime Application Self Protection)} mal implementado: herramientas
    que controlan que los binarios del sistema no cambien durante la ejecución.
    \item \textbf{Verificar Federadores}: protocolos abiertos y antiguos (ej.\ OAuth para
    Google u otras integraciones) que las empresas implementan mal; vulnerabilidades
    conocidas sobre protocolos.
    \item \textbf{Software EOL (End Of Life)}: fuera de soporte, sin actualizaciones de
    seguridad $\rightarrow$ ``muy jugoso'' a la hora de pentestear.
\end{itemize}

\textbf{Resultado del paso 2}: una \textbf{lista de posibles vulnerabilidades} que asumimos
que tiene el sistema.

\subsubsection{Paso 3: Prueba}

\begin{itemize}
    \item \textbf{Priorizar} la lista de posibles vulnerabilidades, según el objetivo de la
    prueba y, por lo general, según el \textbf{nivel de acceso requerido}.
    \item \textbf{Definir cómo probar} la existencia de la vulnerabilidad:
    \begin{itemize}
        \item \textbf{Mejor manera}: analizar documentación o comportamiento.
        \item \textbf{Otra manera}: intentar \textbf{explotarla} $\rightarrow$ último recurso,
        menos eficiente (hay que lograr el exploit como lo haría un atacante).
        \textbf{Excepción}: \emph{Pivoting}.
    \end{itemize}
    \item \textbf{Diseñar el test para que sea lo menos intrusivo posible}: algunas
    vulnerabilidades pueden \textbf{denegar el servicio}.
    \item \textbf{Proceso normal}: resguardar los datos de todo el sistema; documentar los
    requerimientos para detectar la vulnerabilidad; intentar detectarla.
    \item La prueba \textbf{DEBE ser REPETIBLE y CONCLUSIVA}: si funciona una sola vez y no se
    puede replicar, no concluye que haya una vulnerabilidad. Si yo lo puedo repetir,
    naturalmente lo puede repetir un atacante.
\end{itemize}

\begin{destacado}
Hoy por hoy se asume que \textbf{para demostrar una vulnerabilidad hay que tratar de
explotarla} (las ``PoC'' del pentesting), siempre con cuidado, sobre todo en sistemas
productivos: un exploit puede llevar el sistema a un estado inseguro, exponer datos, tirar
abajo la aplicación (atentar contra Confidencialidad, Integridad o Disponibilidad --- el
``CID''). Ejemplos del profe: el caso \textbf{Chernóbil} (era una prueba controlada y
terminó mal) y correr un \texttt{nmap} (escaneo/discovery de red) que tiró abajo la red de
la oficina, afectando la disponibilidad.
\end{destacado}

\begin{ojo}
En pentesting \textbf{NO siempre se intenta explotar} las vulnerabilidades encontradas;
explotar es el \emph{último recurso} y a veces alcanza con analizar documentación o
comportamiento. (V/F, 2015.) En caja negra puede que sí se explote, pero el profe lo matiza:
``intentar explotar es el último recurso, menos eficiente''.
\end{ojo}

\paragraph{Pivoting (excepción).} Es probable que la existencia de una vulnerabilidad provea
más información (ej.\ acceso al S.O.). En ese caso se \textbf{vuelve a la fase de recolección
de información}. \textbf{Pivoting} es el uso de un \textbf{punto de acceso comprometido para
continuar un ataque} (ej.: se compromete el firewall y desde allí se continúa atacando la
red interna). Por eso se encadenan vulnerabilidades y se vuelve al ciclo.

\subsubsection{Paso 4: Generalización}

\begin{itemize}
    \item A medida que las pruebas resultan exitosas, \textbf{emergen patrones}; se ``conectan
    puntos''.
    \item A veces \textbf{dos vulnerabilidades combinadas} constituyen un problema grave.
    Ejemplo: cuenta de invitado habilitada permite conexión remota \emph{+} buffer overflow
    local permite derechos de administrador $\Rightarrow$ desde una cuenta de invitado se
    obtienen permisos de administrador.
    \item Es la \textbf{parte más importante} de la prueba: la ``mano'' del pentester es
    identificar cosas, combinarlas y producir el mayor impacto posible para dar \emph{insight}
    al equipo/empresa que lo contrató.
\end{itemize}

\subsubsection{Paso 5: Eliminación (opcional)}

\begin{itemize}
    \item Por lo general \textbf{solo se incluyen recomendaciones}, porque quien ejecuta la
    prueba no suele ser el diseñador/desarrollador (es el dueño del sistema quien decide qué
    hacer). A veces no se tiene la solución concreta (vulnerabilidades muy nuevas); a veces
    es tan simple como actualizar la versión de una dependencia, a veces es más abstracto
    (ej.\ ``manejá bien la autorización'').
    \item Es importante que queden claros el \textbf{contexto, los detalles y el mecanismo de
    explotación} para: poder corregir el sistema; poder impedir/monitorear mientras tanto;
    verificar si fue explotado; y \textbf{reproducir la prueba} luego de corregir, para
    comprobar que ya no es vulnerable.
    \item Todo termina en un \textbf{reporte de pentesting} con las vulnerabilidades
    encontradas y las recomendaciones.
\end{itemize}

\subsection{Herramientas y técnicas mencionadas}

\begin{itemize}
    \item \textbf{\texttt{nmap}}: programa de Linux que escanea interfaces y puertos, arma una
    topología de la red y hace un \emph{discovery} (qué está abierto/cerrado, qué se conecta
    con qué). \emph{Ojo}: mal usado puede tirar abajo la red.
    \item \textbf{Vulnerability scanning} (ej.\ \texttt{portmapper}) para vulnerabilidades
    conocidas.
    \item \textbf{Ingeniería social} (ver ejemplo de ataque externo).
    \item \textbf{CVE} (catálogo de vulnerabilidades conocidas, con su remediación) y la
    organización \textbf{MITRE}.
    \item \textbf{OWASP Top Ten} (referencia para apps web; Supply Chain como top 3).
\end{itemize}

\subsection{Ejemplos}

\subsubsection{Ejemplo 1: Michigan Terminal System (MTS)}

Sistema operativo que corre en \textbf{mainframes IBM 360/370}. Objetivo de la prueba:
\textbf{obtener acceso a las estructuras de control} que no debería permitirse. Nivel inicial:
\textbf{cuenta autorizada (nivel 3)}. Contaba con la \textbf{aprobación de los administradores}.

\begin{itemize}
    \item \textbf{Paso 1 (Recolección)}: al correr un programa, la memoria se divide en
    segmentos. Segmentos \textbf{0--4}: supervisor, programas de sistema y estado, protegidos
    por hardware. Segmento \textbf{5}: área de trabajo del sistema e información del proceso
    (incluido el nivel de privilegio); el proceso \textbf{no} puede alterarlo. Segmentos
    \textbf{6+}: información del proceso, que el proceso modifica libremente. El segmento 5
    está protegido por memoria virtual: en \textbf{modo sistema (kernel)} el proceso puede
    acceder/modificar y ejecutar funciones privilegiadas; en \textbf{modo usuario} el segmento
    5 no se mapea. Un proceso corre en modo usuario y hace \emph{system calls} (que corren en
    modo sistema). En un system call: se revisan los parámetros (especialmente que los
    punteros pertenezcan a zonas accesibles por el proceso); los parámetros se construyen como
    una \textbf{lista de direcciones} en el segmento de usuario; la lista se pasa como
    dirección en un registro.
    \item \textbf{Paso 2 (Hipótesis)}: ¿qué ocurriría si la dirección de un parámetro
    \textbf{apunta a la lista misma de parámetros}? Sin controles extra, un system call podría
    escribir la dirección del parámetro \textbf{después de haber pasado la validación}. Si se
    puede controlar qué valor escribir, técnicamente se podría leer/escribir cualquier zona del
    \textbf{segmento 5}.
    \item \textbf{Paso 3 (Prueba)}: buscar un system call que use la convención estudiada,
    tenga al menos dos parámetros y altere uno. Se eligió \texttt{line input} (lee una línea de
    texto y retorna número de línea y longitud de línea). \emph{Setup}: hacer que la dirección
    donde se almacena el número de línea sea la de la longitud de línea. \emph{Ejecución}: el
    system call valida los parámetros (todos pertenecen al segmento del proceso); ejecuta
    \texttt{read input line}; la longitud de línea ingresada es el valor a escribir; el número
    de línea se guarda según la lista de parámetros (haciendo que sea una dirección del
    segmento 5, reescribiendo la dirección del otro parámetro), y así se escribe el valor en el
    segmento 5.
    \item \textbf{Paso 4 (Generalización)}: no se puede escribir en los segmentos 0--4
    (protegidos por hardware), pero en el segmento 5 el nivel de privilegio determina si pueden
    hacerse \textbf{llamadas de nivel supervisor} (en lugar de system calls), y una función de
    nivel supervisor \textbf{apaga la protección por hardware}. \textbf{Conclusión}: la
    vulnerabilidad permite a un atacante modificar cualquier dirección de cualquier segmento,
    controlando completamente la computadora.
\end{itemize}

\subsubsection{Ejemplo 2: Ataque externo (ingeniería social)}

\textbf{Objetivo}: determinar si las medidas de seguridad de una empresa eran efectivas para
evitar que un atacante ingrese a un sistema. El test se enfoca en \textbf{políticas y
procedimientos, tanto técnicos como no técnicos}, muy apalancado en \textbf{ingeniería social}
(es más fácil vulnerar a las personas que a los sistemas).

\begin{itemize}
    \item \textbf{Paso 1 (Recolección)}: búsqueda en internet de nombres de empleados y
    directores y teléfono de la sucursal local; del \textbf{reporte anual público} (la empresa
    cotizaba en bolsa) reconstruyeron el organigrama, nombres de proyectos y personas. El tester
    se hizo pasar por nuevo empleado; aprendió que para enviar un documento se necesitaba
    \textbf{número de empleado y centro de costos}. Llamó a la secretaria del director:
    haciéndose pasar por empleado consiguió el número de empleado del director, y haciéndose
    pasar por auditor consiguió el centro de costos. Con eso solicitó enviar el listado a una
    consultora externa.
    \item \textbf{Paso 2 (Hipótesis)}: los empleados nuevos no conocen todos los procesos y
    controles $\rightarrow$ es posible que un nuevo empleado brinde información sensible.
    \item \textbf{Paso 3 (Prueba)}: el tester llamó a \textbf{RRHH} impersonando a la secretaria
    de un director (se quejó de que no le habían informado los nuevos ingresos y los pidió),
    obteniendo los nombres de los nuevos empleados. Luego llamó a esos nuevos empleados haciéndose
    pasar por \textbf{operador del centro de cómputos} y les brindó un ``entrenamiento de
    seguridad'' por teléfono, durante el cual obtuvo: tipos de sistemas usados, número de usuarios,
    \textbf{logins y claves} (incluso induciendo el cambio de una contraseña para conocerla). Se
    probó que los nuevos empleados no estaban al tanto de los procesos/controles.
\end{itemize}

\begin{ojo}
Aunque el objetivo era otro, si se encuentra una falla distinta hay que \textbf{informarla
igual}. Ejemplo del profe: el objetivo era probar la autenticación, pero se encontró que
yendo a una URL del tipo \texttt{/info.txt} desde el browser de cualquier usuario quedaban
expuestos datos de tarjetas $\rightarrow$ se debe reportar más allá del objetivo.
\end{ojo}

\subsection{Limitaciones: el pentesting NO garantiza la seguridad}

\begin{examen}
\faExclam \textbf{GANCHO/TRAMPA de parcial.} El pentesting es una \textbf{buena estrategia
pero NO garantiza la seguridad}.
\prof{Garantizar la seguridad es algo que no se puede hacer jamás.}
\begin{itemize}
    \item Si en un ejercicio alguien afirma que algo \textbf{``garantiza la seguridad''}, es
    incorrecto y \textbf{hay penalización si se les escapa}.
    \item La alternativa con \textbf{más garantía} (aunque \emph{tampoco} total) son las
    \textbf{pruebas / verificación formal}, hoy un área hipercentral por el código generado por
    \textbf{LLMs/agentes}.
    \item No confundir: la verificación formal \emph{prueba la ausencia} de vulnerabilidades
    pero solo en un algoritmo/programa/ambiente acotado y debiendo incluir TODOS los factores;
    el pentesting \emph{nunca} prueba ausencia.
\end{itemize}
\end{examen}

Razones por las que el pentesting no garantiza seguridad (``Para discutir''):
\begin{itemize}
    \item \textbf{No reemplaza} un buen diseño, especificación, implementación ni las pruebas
    (unitarias, de interacción, regresiones). Es una técnica importante para probar el sistema
    \textbf{luego de instalado}; idealmente no sería necesario, en la práctica sí.
    \item Encuentra \textbf{problemas introducidos por la interacción del sistema con los
    usuarios y el ambiente}, que suelen quedar fuera del análisis y pruebas normales.
    \item \textbf{No es sistémico / no es sistemático}: la calidad depende de la
    \textbf{capacidad de los testers para formular hipótesis} (metodología de hipótesis de
    falla). No provee una forma sistemática de revisar un sistema.
    \item \textbf{Cada test es un mundo aparte}: los resultados de un test sirven solo
    marginalmente para otros. Hay herramientas que automatizan \emph{algunos} aspectos, pero
    no todos.
    \item Un pentesting \textbf{vale para ese sistema en ese momento}: si después cambia la API
    o el comportamiento (v2 $\rightarrow$ v3), la mayoría de las pruebas quedan obsoletas.
\end{itemize}

\begin{destacado}
Conclusión del profe: siempre hay que tener buen diseño, respetar los principios de seguridad
y las especificaciones, y hacer buena cobertura de testing; el pentesting se suma \emph{después}
para comprobar que el diseño es bueno. \prof{Siempre algo aparece.}
\end{destacado}

\subsection{Lectura recomendada}

\begin{itemize}
    \item Matt Bishop, \emph{Computer Security: Art and Science}, Capítulo 23 (1--2).
    \item \textbf{OSSTMM} --- Open Source Security Testing Methodology Manual
    (\texttt{http://www.isecom.org/osstmm/}).
\end{itemize}

\end{document}
